\subsubsection{Guiding Principles and Design Patterns}
\label{sec:principles-and-patterns}

The structure described above is not arbitrary; it is intentionally crafted based on a set of industry-standard principles and patterns that ensure high-quality, professional software design. 

\textbf{The SOLID principles}

The SOLID principles are a set of five foundational guidelines that help developers create maintainable and flexible software solutions. This application's architecture adheres to them as follows:

\begin{enumerate} [leftmargin=0.75cm]
    
    \item \textbf{Single Responsibility Principle}
    Each class and layer serves a single purpose and, therefore, has only one reason to change. A clear example is the separation of expense-related responsibilities:

    \begin{itemize} [leftmargin=0.75cm, label=\textbullet]
        
        \item The \texttt{ExpensesController} in the API layer handles HTTP communication with the Presentation layer by receiving requests and returning responses; changes are required only when the API contract evolves.

        \item The \texttt{ExpenseService} in the Application layer is responsible for expense-related business logic; modifications occur only when business rules change.

        \item The \texttt{ExpenseRepository} in the Infrastructure layer manages data persistence; changes are required only when the database or storage service changes.

    \end{itemize}

    \item \textbf{Open/Closed Principle}

    The system is open for extension but closed for modification. Thanks to the interfaces defined in the Application layer, introducing new functionality requires only the implementation of new classes in the Infrastructure layer. The existing, tested code in the Application layer remains unchanged.

    Currently, the \texttt{RepositoryService} in the Application layer depends on the \texttt{IEmailNoti- ficationService}. If, in the future, clients require SMS notifications in addition to email notifications, it is sufficient to introduce a new \texttt{SMSNotificationService} class in the Infrastructure layer that implements a common \texttt{INotificationService} interface. The \texttt{RepositoryService} does not require any modification.

    \item \textbf{Liskov Substitution Principle}

    Any implementation of an interface should be substitutable for another without breaking the application. A clear example is the use of repositories: any class at the Infrastructure layer that implements the \texttt{IExpenseRepository} interface can be used interchangeably by the Application layer. Components that depend on external interfaces are responsible only for knowing how to use them, not for understanding their implementation.

    \item \textbf{Interface Segregation Principle}

    Components should not be forced to depend on unnecessary interfaces. Interfaces are defined at a granular level, such as \texttt{IExpenseRepository} and \texttt{IFileStorage}, rather than combining responsibilities into a single large interface. As a result, components depend only on the interfaces they truly require.

    \item \textbf{Dependency Inversion Principle}

    This is the key principle that holds Clean Architecture together. High-level modules should not depend on low-level modules. For example, the Application layer does not depend on the Infrastructure layer; instead, both depend on abstractions (interfaces). Consequently, as illustrated in Figure \ref{fig:logical-architecture-diagram}, the Infrastructure layer points inward toward the Application layer, resulting in a core business logic that is independent of external concerns.

\end{enumerate}

\textbf{Key Design Patterns}

Several well-known design patterns are used throughout the source code to support the application of these principles:

\begin{itemize} [leftmargin=0.75cm, label=\textbullet]
    
    \item \textbf{Programming to an Interface}

    Components do not depend on concrete classes; instead, they rely on abstract interfaces. These interfaces (e.g., \texttt{IExpenseRepository}) define clear contracts for interacting with external components without exposing services' internal logic. This approach is the primary mechanism for achieving a decoupled system. 

    \item \textbf{Service Pattern}

    The logic in the Application layer is organized into services (e.g., \texttt{ExpenseService}). Each service groups related business operations, providing a clear and easy-to-understand structure. 

    \item \textbf{DTO Pattern}

    Data Transfer Objects (DTOs) are simple data containers that protect internal domain models from external exposure. Their primary purpose is to facilitate data transfer between the API and Application layers, establishing a clear and consistent contract between them. 

    \item \textbf{Domain-Driven Design (DDD)}

    The Domain layer is structured according to Domain-Driven Design principles. It encapsulates the core business logic within a pure, isolated layer.

    \item \textbf{Repository Pattern}

    Repositories implement interfaces such as \texttt{IExpenseRepository} and encapsulate the logic behind these abstractions. For example, the Application layer can request to save an expense via \texttt{IExpenseRepository} without needing to know which database or storage mechanism is used. 

    \item \textbf{Dependency Injection (DI)} 

    Dependency Injection enables components to depend on external sources (e.g., \texttt{Expen- seService} depends on a repository for data operations). It is the practical implementation of the Dependency Inversion Principle, allowing layers to remain decoupled in the code while being connected and functional at runtime. This approach clarifies component responsibilities, reduces the need for code changes, and simplifies testing by removing dependencies on internal implementations. 

\end{itemize}

All these patterns collectively support and reinforce the adoption of Clean Architecture principles. For a complete overview of the system structure applying these principles and patterns, refer to the UML Class Diagram in Section \ref{sec:backend-class-diagram}.
